\section{Deciding the forbidden-family question}

The comparability-graph correspondence reduces the question to a hereditary
class of graphs. A graph can have many binary union/join representations,
and those representations need not embed into a common tree even when their
graphs have a common host. We therefore decompose a language of representations
into directed pieces and compare their \emph{graph images}. This section
gives the effective steps, including the passage between different binary
representations of the same graph.

We use two standard effective results. First, for a regular language of
ranked trees, one can compute its homeomorphic downward closure as a finite
union of regular ideals. The precise input is the construction of
Goubault-Larrecq and Schmitz~\cite{GLS}, Section~5: Proposition~18 constructs
simple tree regular expressions; Fact~19 and Lemma~20 give a terminating
normalization; Theorem~23 identifies the resulting tree products as ideals.
Second, the Parikh image of a context-free language has an effective
Presburger description. One explicit construction is given by Esparza,
Ganty, Kiefer and Luttenberger~\cite{Parikh}, Theorem~1.1 and Section~3.
Presburger arithmetic is first-order arithmetic with addition and order on
the natural numbers; its sentences are decidable. These inputs concern tree
languages and finite grammars. The graph-containment and JEP transfers below
are proved here.

\subsection{Binary representations and ideals}
Use the ranked alphabet $\{a:0,U:2,J:2\}$. Each occurrence of $a$ represents
one vertex. The operations $U$ and $J$ represent disjoint union and join,
respectively. Write $\E(t)$ for the graph represented by a term $t$.
Children are ordered in a term, although graph evaluation forgets their
order. The empty graph is treated separately.

The ordered homeomorphic relation $\preceq$ on these terms has the following
recursive definition. A source can embed into either child of a target;
alternatively, equal root symbols can be coupled with childwise embeddings.
Thus $a\preceq a$, and for $o\in\{U,J\}$ the coupling clause is
\[
 o(s_1,s_2)\preceq o(t_1,t_2)
 \quad\hbox{if }s_1\preceq t_1\hbox{ and }s_2\preceq t_2.
\]
This is exactly the ranked relation in Example~7 of~\cite{GLS}.
It preserves the root operation at each coupled node. An \emph{ideal}
is a nonempty downward-closed directed set under the specified relation.

\begin{lemma}\label{lem:parse-pruning}
If $s\preceq t$, then $\E(s)$ embeds as an induced subgraph of $\E(t)$.
For every nonempty subset $B$ of the leaves of $t$, there is a term
$t|B\preceq t$ representing exactly $\E(t)[B]$.
\end{lemma}
\begin{proof}
For the first assertion, diving includes a child graph as an induced
subgraph of its parent. Coupling combines two child embeddings on disjoint
vertex sets: the common operation puts no edges, or all edges, between
those sets on both sides. Induction proves the assertion.

For the second, delete leaves outside $B$ recursively. At a node discard
empty child results; return the surviving child if there is only one,
and retain the original binary operation if both survive. Suppressing a
unary node keeps the surviving child's own operation. Restriction commutes
with both graph operations, which proves the graph identity. The dive and
coupling clauses prove $t|B\preceq t$ by induction. Adjacent equal
operations need not be flattened: the result is a binary representation.
\end{proof}

Let $\mathcal F$ be a finite forbidden graph family containing $P_4,C_4$,
and put $\mathcal K=\Forb(\mathcal F)$. Every nonempty member of
$\mathcal K$ has a binary representation: the construction in
Lemma~\ref{lem:bridge} recursively takes disjoint unions and adds a
universal vertex by joining with $a$.

\begin{lemma}\label{lem:parse-regular}
The language $L=\{t:\E(t)\in\mathcal K\}$ is effectively regular and
downward closed under $\preceq$.
\end{lemma}
\begin{proof}
Choose $n\geq1$ at least as large as every forbidden graph. A state records
the set of induced graph isomorphism types with at most $n$ vertices present
in the represented graph, including the empty type. There are finitely many
such types, and canonical representatives and isomorphism tests are finite
computations. The leaf state is $\{\varnothing,K_1\}$. For either operation
$o$, the transition on states $S,T$ is
\[
 \{o(G,H):G\in S,\ H\in T,\ |G|+|H|\leq n\},
\]
with graph types identified up to isomorphism. The empty graph is an
identity for both operations. Every induced restriction of an operation
splits into restrictions of its two arguments, so these transitions compute
the stated profile exactly. Accept states containing no member of
$\mathcal F$. This constructs the required finite automaton. Heredity
and Lemma~\ref{lem:parse-pruning} give downward closure.
\end{proof}

Apply the effective ideal-decomposition theorem to obtain
\[
                     L=I_1\cup\cdots\cup I_r,                 \tag{1}
\]
where the $I_i$ are regular term ideals with effectively available automata.
If a forbidden graph is empty, $\mathcal K$ is empty and has vacuous JEP.
Otherwise it contains the empty graph. If $L$ is empty, $\mathcal K$
contains only that graph and has JEP too. These cases are decidable;
henceforth $r\geq1$.

\begin{lemma}\label{lem:image-cover}
Put $D_i=\{\varnothing\}\cup\{\E(t):t\in I_i\}$, on isomorphism types.
Each $D_i$ is a hereditary graph class with JEP, their union is
$\mathcal K$, and
\[
 \mathcal K\text{ has JEP}\quad\Longleftrightarrow\quad
                     \mathcal K=D_i\text{ for some }i.        \tag{2}
\]
\end{lemma}
\begin{proof}
Containment and covering follow from (1) and the representation of every
nonempty graph in $\mathcal K$. For heredity take a representation in
$I_i$ and apply the pruning part of Lemma~\ref{lem:parse-pruning}; the
result stays in $I_i$. For JEP take two representing terms in $I_i$,
use its directedness, and evaluate a common term host. Empty graphs cause
no difficulty.

Equality with one $D_i$ implies JEP. Conversely, suppose $\mathcal K$
has JEP but each $D_i$ is proper. Choose $G_i\in\mathcal K\setminus D_i$.
Iterating JEP supplies one $H\in\mathcal K$ containing all the finitely
many $G_i$. The cover puts $H$ in some $D_j$; heredity then puts $G_j$
there, a contradiction.
\end{proof}

The remaining task is to decide the finitely many inclusions
$\mathcal K\subseteq D_i$. Individual binary representations cannot be
used for this test, since a graph may have another representation accepted
by $I_i$. We now give the required exact test.

\subsection{Recognizing graph images of regular term languages}
A \emph{reduced cotree} has leaves $a$ and unordered internal nodes
$U,J$ of arity at least two, with adjacent internal operations unequal.
Flattening adjacent equal operations turns a binary representation into
such a cotree without changing its graph.

\begin{lemma}\label{lem:canonical-cotree}
Every graph with a union/join representation has a unique reduced cotree
up to isomorphism. All its binary representations are obtained by
independently representing each cotree child, then choosing a binary
association and a permutation of the children at each operation node.
\end{lemma}
\begin{proof}
At a reduced $U$ root each child graph is connected: a child is a leaf or
has a join root. Its children are therefore precisely the graph's connected
components. At a $J$ root the graph is connected and its children are
precisely the connected components of the complement. These decompositions
recover the root operation and its child vertex sets from the graph.
Induction recovers every child cotree; zero and one vertices are immediate.
This proves uniqueness. Flattening any binary representation yields this
cotree, and reversing each flattening uses precisely binary association
and permutation within a single operation. This proves the last assertion.
\end{proof}

Let $M$ be a deterministic complete binary tree automaton with finite
state set $Q$, leaf state $q_a$, transitions $\delta_U,\delta_J$, and
accepting set $A_M$. For a reduced cotree $T$ define
\[
 S_M(T)=\{q\in Q:\text{some binary representation of }T
                         \text{ has state }q\}.              \tag{3}
\]
Storing the full subset accounts for all possible representations.

\begin{lemma}\label{lem:parikh-transition}
For each $o\in\{U,J\}$ and $q\in Q$ one can compute a Presburger
predicate $P_{o,q}$ on the counts of child state subsets such that, if
an $o$ node has children $T_1,\ldots,T_k$, then
\[
 q\in S_M(T)\quad\Longleftrightarrow\quad
 P_{o,q}\bigl((c_B)_{B\subseteq Q}\bigr),\qquad
 c_B=|\{j:S_M(T_j)=B\}|.                                   \tag{4}
\]
Equal child occurrences are counted with their full multiplicity.
\end{lemma}
\begin{proof}
Use nonterminals $N_q$ for $q\in Q$ and one terminal $b_B$ for every
subset $B\subseteq Q$. For the operation $o$ form the finite grammar
\[
 N_{\delta_o(r,s)}\longrightarrow N_rN_s\quad(r,s\in Q),
 \qquad N_q\longrightarrow b_B\quad(q\in B).
\]
Compute the Parikh predicate for start nonterminal $N_q$. A binary
representation of the node consists of a binary $o$ skeleton, with each
child represented at a chosen state belonging to its subset. Reading
that skeleton gives a derivation with the stated terminal counts.

Conversely a derivation with those counts supplies an $o$ skeleton and
one selected state for each terminal occurrence. Match occurrences to
the actual children of the same subset type. For each occurrence choose
a child representation attaining its selected state, and assemble the
representations in the skeleton's order. This is a representation of
$T$ with root state $q$. Lemma~\ref{lem:canonical-cotree} ensures that
these are exactly the allowed representations. The effective Parikh
theorem therefore proves both the correspondence and effectivity.
\end{proof}

The resulting deterministic cotree automaton has states $(\ell,S)$,
where $\ell\in\{a,U,J\}$ records the root kind and $S\subseteq Q$.
Its leaf state is $(a,\{q_a\})$. At an $o$ node, transition to a
specified subset $S$ is expressed by
\[
 \bigwedge_{q\in S}P_{o,q}(c)\ \land\
 \bigwedge_{q\notin S}\neg P_{o,q}(c).                        \tag{5}
\]
Require at least two children and exclude children with operation $o$.
The counts $c_B$ are sums of the counts for the permitted child root
kinds with subset $B$. These are linear expressions, so the horizontal
transitions remain Presburger formulas. Accept when $S\cap A_M$ is
nonempty. Induction using (4) shows that this automaton accepts exactly
the graphs represented by terms accepted by $M$, in their canonical
cotree form.

\begin{lemma}\label{lem:presburger-inclusion}
Inclusion of languages recognized by these deterministic cotree
automata is decidable.
\end{lemma}
\begin{proof}
Work throughout on the domain of reduced cotrees. Complement changes
the accepting set. For an intersection, use pairs of states; replace
each component child-state count by the sum of pair-state counts
projecting to it, and conjoin the transition formulas. Root kinds,
arity and alternation agree in the product. These are finite
Presburger constructions.

For emptiness, start with the actual leaf states and compute a least
fixed point of productive states. Given the current finite set $R$,
add an internal state if its transition formula has a nonnegative
count vector of total at least two, with count zero on every child
type outside $R$. Presburger decision settles each such question.
There are finitely many states, so the process terminates.

Each newly added state has a finite cotree representative: take the
required number of copies of representatives for its already productive
child states. Conversely induction on cotree height puts every state
of an actual finite cotree into the fixed point. Thus a language is
nonempty exactly when an accepting state is productive. Apply this to
the first language intersected with the complement of the second.
\end{proof}

\subsection{The resulting algorithm}
For the automata of $L$ and each $I_i$, construct the cotree automata
(3)--(5). They recognize $\mathcal K$ and $D_i$, respectively, with
the common empty graph omitted. Lemma~\ref{lem:presburger-inclusion}
decides $\mathcal K\subseteq D_i$. Since $D_i\subseteq\mathcal K$
already, this tests equality. By (2), return yes if some equality
holds and no otherwise. Every stage terminates, by the cited effective
theorems or the explicit finite-state procedure just proved.

For example, the independent graphs with at most four vertices have
JEP. Their binary representations include a four-leaf balanced tree and
a four-leaf comb, which cannot have a common binary host with at most
four leaves under $\preceq$. Our algorithm decomposes the representation
language instead of testing its JEP. In this finite example the maximal
binary representations generate principal term ideals. Each such ideal
has graph image exactly the independent graphs of sizes one through
four, by pruning and leaf counts. Their graph images therefore give
the correct affirmative verdict.

\begin{theorem}\label{thm:decision}
There is an algorithm which, given any finite family $\mathcal U$ of
finite order forests, decides whether $\mathcal C(\mathcal U)$ has JEP.
More precisely, set
\[
 \mathcal F=\{P_4,C_4\}\cup\{\Comp(U):U\in\mathcal U\}.
\]
Then $\mathcal C(\mathcal U)$ has JEP if and only if $\Forb(\mathcal F)$
has JEP.
\end{theorem}
\begin{proof}
Lemma~\ref{lem:bridge} gives an effective correspondence on isomorphism
types between finite forests and $\Forb(P_4,C_4)$, preserving and
reflecting existence of embeddings. Thus avoiding $\mathcal U$
corresponds exactly to avoiding the displayed graph family, and common
hosts correspond in both directions. Compute this finite graph list
and run the terminating graph procedure established above.
\end{proof}

An input family need not be an antichain: remove redundant forbidden
members by finite embedding tests. The algorithm supplies an effective
criterion, with no claim that its large intermediate automata are
practical or that the decision procedure has been implemented.
